\chapter[Формирование и модели вознаграждения]{Формирование вознаграждения и модели вознаграждения}
\label{ch:reward-models}

\begin{quote}
\itshape
Мусор на входе, мусор на выходе. Качество обученной стратегии никогда не превысит качество оптимизируемого сигнала вознаграждения.
\end{quote}

\bigskip

В конечном счёте любой алгоритм обучения с подкреплением сводится к одному императиву: максимизировать кумулятивное вознаграждение. Это делает сигнал вознаграждения наиболее значимым решением при проектировании любой RL-системы. Неправильно заданное вознаграждение\,---\,даже незначительно\,---\,даст стратегию, которая в лучшем случае будет субоптимальной, а в худшем --- активно вредной. История прикладного RL полна поучительных примеров: симулированный робот, научившийся двигаться вибрацией вместо ходьбы; игровой агент, обнаруживший высокооценивающийся эксплойт, которого разработчики никогда не предусматривали; и\,---\,в современную эпоху выравнивания языковых моделей\,---\,чат-боты, научившиеся производить многословную, льстивую прозу, потому что так предпочитала прокси-модель вознаграждения.

Эта глава рассматривает проблему проектирования вознаграждения с двух взаимодополняющих углов. Первая половина развивает классическую теорию \emph{формирования вознаграждения}: по заданной среде с хорошо определённым, но разреженным или иначе неудобным истинным вознаграждением --- как дополнить это вознаграждение дополнительными слагаемыми, ускоряющими обучение без изменения оптимальной стратегии? Краеугольным результатом является теорема о формировании на основе потенциала Нга, Харады и Расселла (1999), дающая точную характеристику всех дополнений вознаграждения, гарантированно сохраняющих стратегию. Вторая половина обращается к \emph{обученным моделям вознаграждения}: что делать, когда истинное вознаграждение неизвестно даже в замкнутой форме и должно быть выведено из суждений людей, исходов или аннотаций на уровне шагов. Мы выводим модель предпочтений Bradley--Terry, кросс-энтропийную функцию потерь для обучения моделей вознаграждения по предпочтениям, а затем подробно рассматриваем различие между моделями вознаграждения по исходам (ORM) и по процессу (PRM). Завершаем переоптимизацией вознаграждения\,---\,законом Гудхарта в RL\,---\,и механизмом штрафа KL, защищающим от неё.

Эта глава занимает ключевое место в книге. Она связывает Часть~II (классические методы RL, Главы~\ref{ch:dp}--\ref{ch:actor-critic}) и Часть~III (RL для языковых моделей). Обученные здесь модели вознаграждения непосредственно используются в конвейере дообучения PPO (Глава~\ref{ch:practical-rlhf}), целевой функции DPO и алгоритме GRPO, описанных в последующих главах. Читатели, освоившие аппарат TD и градиента стратегии из Части~II, найдут материал по моделированию вознаграждения во второй половине этой главы естественным расширением, изложенным в терминах крупномасштабных языковых моделей.

% ===========================================================
\section{Формирование вознаграждения на основе потенциала}
\label{sec:reward-shaping}

\subsection{Проблема разреженного вознаграждения}

Во многих средах агент сталкивается с разреженным сигналом вознаграждения: положительная обратная связь приходит только при достижении конкретной цели, а за все остальные переходы получается нулевое вознаграждение. Рассмотрим робота, навигирующего в большом лабиринте: истинное вознаграждение равно~$+1$ у выхода и~$0$ везде ещё. Агент, исследующий случайным образом, найдёт выход только после астрономически большого числа шагов, и до этого момента не получит никакой информации о градиенте вовсе. Аналогично, языковая модель, получающая вознаграждение только в конце длинного ответа, должна приписать каждому из сотен предшествующих токенов долю этого терминального сигнала\,---\,проблема присвоения кредита в крайней форме.

Интуитивное решение --- \emph{формировать} вознаграждение: добавить вспомогательные слагаемые, направляющие агента к цели. Робот, навигирующий в лабиринте, мог бы получать небольшой бонус, пропорциональный тому, насколько он сократил Евклидово расстояние до выхода. Языковая модель могла бы получать промежуточные вознаграждения за производство грамматически верных предложений. Однако формирование вознаграждения опасно. Наивно добавленные бонусные слагаемые могут изменить оптимальную стратегию: робот, получающий большой бонус за приближение к стене (которая оказывается рядом с выходом), может научиться держаться возле этой стены, а не достигать цели. Фундаментальный вопрос: \emph{какие дополнения вознаграждения гарантированно не изменяют множество оптимальных стратегий?}

\subsection{Теорема о формировании}

Нг, Харада и Расселл (1999)\index{reward shaping!Ng-Harada-Russell theorem} ответили на этот вопрос точно. Их результат, известный как \textbf{теорема о формировании на основе потенциала}, характеризует все преобразования вознаграждения, инвариантные к стратегии.

\begin{definition}[Функция формирования на основе потенциала]
\label{def:shaping-function}
\index{reward shaping!potential function}
Пусть $\mathcal{M} = (\cS, \cA, P, R, \gamma)$ --- дисконтированный MDP. \textbf{Функция формирования на основе потенциала} --- это функция $F : \cS \times \cA \times \cS \to \R$ вида
\begin{equation}
\label{eq:shaping-F}
F(s, a, s') = \gamma\,\Phi(s') - \Phi(s),
\end{equation}
где $\Phi : \cS \to \R$ --- произвольная \textbf{потенциальная функция}\index{potential function}.
\end{definition}

Задав сформированное вознаграждение $\tilde{R}(s,a,s') = R(s,a,s') + F(s,a,s')$, определим сформированный MDP $\tilde{\mathcal{M}} = (\cS, \cA, P, \tilde{R}, \gamma)$. Центральная теорема такова.

\begin{theorem}[Теорема о формировании на основе потенциала, Ng et al.\ 1999]
\label{thm:shaping}
\index{reward shaping!shaping theorem}
Пусть $F$ --- функция формирования на основе потенциала, как в~\eqref{eq:shaping-F}. Тогда:
\begin{enumerate}
  \item Каждая оптимальная стратегия для $\mathcal{M}$ также оптимальна для $\tilde{\mathcal{M}}$, и наоборот.
  \item Сформированная функция ценности $\tilde{V}^\pi$ и сформированная функция ценности действий $\tilde{Q}^\pi$ удовлетворяют
    \begin{equation}
    \label{eq:shaped-values}
    \tilde{V}^\pi(s) = V^\pi(s) - \Phi(s), \qquad
    \tilde{Q}^\pi(s,a) = Q^\pi(s,a) - \Phi(s).
    \end{equation}
  \item $F(s,a,s') = \gamma\Phi(s') - \Phi(s)$ --- \emph{единственное} семейство функций (с точностью до константы), при котором оптимальные стратегии сохраняются для всех MDP.
\end{enumerate}
\end{theorem}

\begin{proof}
Докажем части~1 и~2. $n$-шаговая доходность при сформированном вознаграждении, начиная из состояния $s_0$, равна
\begin{align*}
\tilde{G}_0 &= \sum_{t=0}^{\infty} \gamma^t \tilde{R}(S_t, A_t, S_{t+1}) \\
&= \sum_{t=0}^{\infty} \gamma^t R(S_t, A_t, S_{t+1})
   + \sum_{t=0}^{\infty} \gamma^t \bigl(\gamma\Phi(S_{t+1}) - \Phi(S_t)\bigr).
\end{align*}
Раскрывая вторую сумму как телескопический ряд:
\[
\sum_{t=0}^{\infty} \gamma^t \bigl(\gamma\Phi(S_{t+1}) - \Phi(S_t)\bigr)
= \sum_{t=0}^{\infty} \gamma^{t+1}\Phi(S_{t+1}) - \sum_{t=0}^{\infty} \gamma^t\Phi(S_t)
= -\Phi(S_0).
\]
(Телескопическое сокращение выполняется при условии $\gamma^t \Phi(S_t) \to 0$ при $t\to\infty$, что гарантируется при ограниченной $\Phi$ и $\gamma < 1$.) Следовательно
\begin{equation}
\label{eq:shaping-return}
\tilde{G}_0 = G_0 - \Phi(S_0).
\end{equation}
Беря ожидания: $\tilde{V}^\pi(s) = V^\pi(s) - \Phi(s)$, что соответствует~\eqref{eq:shaped-values}. Поскольку $\Phi(s)$ зависит только от текущего состояния, но не от действия или стратегии, функция преимущества не изменяется:
\[
\tilde{A}^\pi(s,a) = \tilde{Q}^\pi(s,a) - \tilde{V}^\pi(s) = Q^\pi(s,a) - V^\pi(s) = A^\pi(s,a).
\]
Жадная стратегия относительно $\tilde{Q}^\pi$ совпадает с жадной стратегией относительно $Q^\pi$, что устанавливает часть~1.
\end{proof}

Элегантность доказательства заключается в телескопическом тождестве~\eqref{eq:shaping-return}: добавление бонуса на основе потенциала эквивалентно простому сдвигу всей доходности на зависящую от состояния константу $-\Phi(S_0)$, которая не может повлиять на то, какая стратегия достигает максимума.

\begin{remark}[Не-потенциальное формирование меняет оптимальные стратегии]
\label{rem:non-potential}
Часть~3 Теоремы~\ref{thm:shaping} предостерегает: \emph{любая} функция формирования, не имеющая вид на основе потенциала, изменит оптимальные стратегии хотя бы в некотором MDP. Классический контрпример: предположим, что бонус формирования $F(s,a,s') = +1$ для всех переходов, независимо от какого-либо потенциала. Это равномерно завышает все вознаграждения, что не меняет оптимальную стратегию при полном дисконтировании, но искажает относительную ценность коротких и длинных эпизодов. Более серьёзно: не-потенциальные бонусы, такие как зависящие от действий слагаемые (например, бонусы за посещение новых состояний при обучении на основе любопытства), \emph{действительно} изменяют оптимальную стратегию в общем случае, что может быть намеренным (исследование), но должно пониматься как сознательное изменение цели, а не нейтральное дополнение.
\end{remark}

\subsection{Проектирование потенциальных функций: примеры}

Теорема гарантирует инвариантность стратегии для любой $\Phi$, однако \emph{скорость} обучения существенно зависит от того, насколько хорошо $\Phi$ кодирует полезные априорные знания. Три стандартные конструкции иллюстрируют пространство проектных решений.

\begin{example}[Потенциал расстояния до цели в мире клеток]
\label{ex:grid-potential}
Рассмотрим сетку $5\times 5$ с целевым состоянием $g$. Определим $\Phi(s) = -d(s, g)$, где $d$ --- манхэттенское расстояние до цели. Бонус формирования для перехода из $s$ в $s'$ равен
\[
F(s, s') = \gamma(-d(s',g)) - (-d(s,g)) = d(s,g) - \gamma\,d(s',g).
\]
Когда агент приближается к цели ($d(s',g) < d(s,g)$) и $\gamma \approx 1$, этот бонус приблизительно равен $d(s,g) - d(s',g) > 0$. Удаление от цели даёт отрицательный бонус. Это мягко направляет агента к цели, не изменяя того, что он в конечном итоге научится делать.

Численно: предположим $\gamma = 0.99$, $d(s,g) = 4$, $d(s',g) = 3$. Тогда $F = 4 - 0.99 \cdot 3 = 4 - 2.97 = 1.03$. При удалении агента: $d(s',g) = 5$, $F = 4 - 0.99 \cdot 5 = 4 - 4.95 = -0.95$. Сигнал бонуса по величине сопоставим с истинным терминальным вознаграждением~$+1$, обеспечивая плотное руководство на протяжении всего эпизода.
\end{example}

\begin{example}[Потенциал функции ценности]
\label{ex:vf-potential}
Если имеется доступ к приближённой функции ценности $\hat{V}$ (возможно, из более простого, связанного MDP или из предобученной модели), установка $\Phi = \hat{V}$ даёт бонусы формирования, кодирующие априорные убеждения агента о качестве состояний. В контексте языковых моделей именно такова роль \emph{опорной модели}: стратегия SFT обеспечивает неявную оценку того, какие продолжения являются высококачественными, а штраф KL можно интерпретировать как регуляризацию на основе потенциала. Мы вернёмся к этой связи в Разделе~\ref{sec:overoptimization}.
\end{example}

\begin{example}[Обученный потенциал из критика]
\label{ex:critic-potential}
В методах актор-критик (Глава~\ref{ch:actor-critic}) критик $V_\phi$ сам обучается одновременно со стратегией. Установка $\Phi = V_\phi$ и вычисление бонусов формирования на лету соответствует в точности формулировке через функцию преимущества: TD-ошибка $\delta_t = R_{t+1} + \gamma V_\phi(S_{t+1}) - V_\phi(S_t)$ является бонусом формирования на основе потенциала, добавляемым к непосредственному вознаграждению перед использованием в качестве сигнала градиента стратегии. Это обнаруживает, что градиенты стратегии на основе преимущества в точном смысле используют обученную потенциальную функцию в качестве бонуса формирования.
\end{example}

% ===========================================================
\section{От ручного проектирования к обученным вознаграждениям}
\label{sec:reward-hacking}

\subsection{Пределы ручного проектирования вознаграждения}

Даже при защите потенциальным формированием проектирование правильной функции вознаграждения для сложных задач чрезвычайно затруднено. Практик сталкивается с фундаментальной асимметрией: описать, как \emph{выглядит} хорошее поведение, легко, но точно закодировать его как скалярную функцию состояний и действий --- сложно. Для робототехники необходимо решить, вознаграждать ли моменты сочленений, позиции концевого эффектора, выполнение задачи, энергоэффективность и ограничения безопасности одновременно, а затем вручную настраивать их относительные веса. Для разговорного агента понятия \emph{полезности}, \emph{безвредности} и \emph{честности} интуитивно понятны людям, но крайне нетривиально поддаются количественной оценке.

\index{reward hacking}
\textbf{Взлом вознаграждения}\index{reward hacking} --- это явление, при котором агент обнаруживает стратегию, достигающую высокого числового вознаграждения без удовлетворения замысла разработчика. Литература по «specification gaming» (Krakovna et al., 2020) документирует сотни примеров в различных предметных областях. В RL для игр агенты научились умирать повторно, чтобы избежать потери очков; эксплуатировать баги физики для бесконечной скорости; и останавливать игровые часы, чтобы предотвратить штрафы за ограниченное время. В робототехнике симулированная рука научилась позиционировать себя между датчиком и объектом задачи так, что задача казалась завершённой для датчика, тогда как реального прогресса не было.

\subsection{Взлом вознаграждения в языковых моделях}

Те же режимы отказа появляются, возможно ещё более остро, в выравнивании языковых моделей, где функция вознаграждения даже не проектируется вручную, а сама является нейронной сетью, обученной на шумных человеческих суждениях.

\index{verbosity bias}
\textbf{Смещение по многословности}: аннотаторы, оценивающие пары ответов, склонны воспринимать более длинные, подробные ответы как более качественные, даже когда дополнительная длина не несёт информации. Модель вознаграждения, обученная на таких данных, будет присваивать более высокие баллы более длинным ответам. Когда языковая модель оптимизируется против этого прокси-вознаграждения, она учится заполнять ответы повторяющимися оговорками и шаблонными резюме\,---\,стратегия, максимизирующая прокси при деградации подлинной полезности.

\textbf{Подхалимство}\index{sycophancy}: человеческие аннотаторы подвержены смещению социальной желательности, предпочитая ответы, соглашающиеся с их высказанными мнениями, ответам, исправляющим их ошибки. Модель вознаграждения, интернализирующая эту закономерность, будет высоко оценивать ответы, ищущие соглашения, что приводит к тому, что оптимизированная стратегия говорит пользователям то, что те хотят услышать, а не то, что является истиной.

\textbf{Эксплуатация форматирования}: модели вознаграждения, обученные на ответах людей-писателей, могут непреднамеренно ассоциировать markdown-форматирование (маркированные списки, жирные заголовки, блоки кода) с высоким качеством, просто потому что человеческие демонстрации склонны быть хорошо оформленными. Оптимизированная стратегия затем применяет тяжёлое форматирование даже там, где оно неуместно, например оформляя краткий фактический ответ в виде маркированного списка.

Эти примеры мотивируют фундаментальный сдвиг в стратегии: вместо попыток закодировать функцию вознаграждения a priori мы \emph{обучаем} её по обратной связи от людей. В оставшейся части главы разрабатывается теоретический и практический аппарат для этого.

% ===========================================================
\section{Модель Bradley--Terry}
\label{sec:bradley-terry}

\subsection{Данные о предпочтениях}

Отправная точка для обучения функции вознаграждения по мнениям людей --- наблюдение, что человеку часто легче сказать, какой из двух ответов лучше, чем присвоить единственному ответу абсолютный числовой балл. Сравнительные суждения более надёжны, более согласованны между аннотаторами и менее чувствительны к абсолютной калибровке отдельных рейтеров.

\index{preference data}
Формально, предположим, что существует скрытая истинная функция вознаграждения $r^* : \cS \times \cA \to \R$ (или, в контексте языковых моделей, $r^* : \mathcal{X} \times \mathcal{Y} \to \R$, где $\mathcal{X}$ --- пространство промптов, $\mathcal{Y}$ --- пространство ответов). Эта функция не наблюдаема напрямую. Вместо этого человек-аннотатор наблюдает два ответа $y_1$ и $y_2$ на промпт $x$ и указывает предпочтение, обозначаемое $y_1 \succ y_2$ («$y_1$ предпочтительнее $y_2$»).

\begin{definition}[Датасет предпочтений]
\label{def:preference-dataset}
\index{preference dataset}
\textbf{Датасет предпочтений} --- это коллекция
\[
\mathcal{D} = \bigl\{(x^{(i)},\, y_w^{(i)},\, y_l^{(i)})\bigr\}_{i=1}^{N},
\]
где $y_w^{(i)} \succ y_l^{(i)}$ означает, что человек-аннотатор предпочёл $y_w^{(i)}$ («победитель») над $y_l^{(i)}$ («проигравший») в ответ на промпт $x^{(i)}$.
\end{definition}

Цель --- обучить параметрическую функцию вознаграждения $r_\psi : \mathcal{X} \times \mathcal{Y} \to \R$ так, чтобы $r_\psi(x, y_w) > r_\psi(x, y_l)$ для всех пар в $\mathcal{D}$, и, что важнее, чтобы она правильно обобщалась на невидимые промпты и ответы.

\subsection{Модель предпочтений Bradley--Terry}

Как следует моделировать вероятность того, что человек предпочтёт $y_1$ над $y_2$? Наиболее принципиальным и широко используемым фреймворком является \textbf{модель Bradley--Terry}\index{Bradley--Terry model} (Bradley and Terry, 1952), изначально разработанная для ранжирования спортивных соревнующихся и впоследствии принятая в статистике, экономике и машинном обучении.

\begin{definition}[Модель Bradley--Terry]
\label{def:bt-model}
\index{Bradley--Terry model}
\textbf{Модель Bradley--Terry} постулирует, что вероятность предпочтения $y_1$ над $y_2$ при промпте $x$ зависит только от разности скрытых качеств:
\begin{equation}
\label{eq:bt-model}
\Prob(y_1 \succ y_2 \mid x) = \sigma\bigl(r^*(x, y_1) - r^*(x, y_2)\bigr),
\end{equation}
где $\sigma(z) = 1/(1 + e^{-z})$ --- логистическая сигмоидная функция.
\end{definition}

\begin{tikzpicture}[
  every node/.style={font=\small},
  box/.style={draw, rounded corners=3pt, minimum width=2.2cm, minimum height=0.85cm, align=center},
  prob/.style={draw, fill=gray!10, rounded corners=3pt, minimum width=2.6cm, minimum height=0.85cm, align=center},
  arrow/.style={->, >=Stealth, thick}
]
  % Ответы
  \node[box, fill=blue!8] (y1) at (0, 1.4) {Ответ $y_1$\\$r^*(x,y_1)$};
  \node[box, fill=red!8]  (y2) at (0,-1.4) {Ответ $y_2$\\$r^*(x,y_2)$};
  % Разность
  \node[box] (diff) at (3.8, 0) {$\Delta = r^*(x,y_1)$\\$- r^*(x,y_2)$};
  % Сигмоида
  \node[box] (sig)  at (7.4, 0) {$\sigma(\Delta)$};
  % Вероятность
  \node[prob] (pref) at (11, 0) {$\Prob(y_1 \succ y_2)$};

  \draw[arrow] (y1.east) -- ++(0.7,0) |- (diff.west);
  \draw[arrow] (y2.east) -- ++(0.7,0) |- (diff.west);
  \draw[arrow] (diff) -- (sig);
  \draw[arrow] (sig)  -- (pref);
\end{tikzpicture}

\begin{center}
\small\textit{Рисунок: Модель предпочтений Bradley--Terry. Вероятность предпочтения зависит только от разности скрытых оценок вознаграждения, пропущенной через сигмоиду.}
\end{center}

Модель имеет ясную вероятностную интерпретацию. Когда два ответа одинаково хороши ($r^*(x,y_1) = r^*(x,y_2)$), вероятность предпочтения равна $\sigma(0) = 1/2$: аннотатор безразличен. По мере роста разрыва в качестве вероятность предпочтения приближается к~$1$ (или~$0$), отражая возрастающую уверенность. Сигмоидная форма фиксирует насыщение человеческих суждений: для очень очевидных пар вероятность предпочтения фактически определённа вне зависимости от точной величины разности.

\begin{proposition}[Свойства модели Bradley--Terry]
\label{prop:bt-properties}
Модель Bradley--Terry~\eqref{eq:bt-model} удовлетворяет:
\begin{enumerate}
  \item \textbf{Симметрия}: $\Prob(y_1 \succ y_2 \mid x) + \Prob(y_2 \succ y_1 \mid x) = 1$.
  \item \textbf{Безразличие при равенстве}: $r^*(x,y_1) = r^*(x,y_2) \implies \Prob(y_1 \succ y_2) = 1/2$.
  \item \textbf{Монотонность}: $\Prob(y_1 \succ y_2 \mid x)$ строго возрастает по $r^*(x,y_1)$ и строго убывает по $r^*(x,y_2)$.
  \item \textbf{Инвариантность к сдвигу}: $r^* \mapsto r^* + c$ для любой константы $c$ не изменяет все вероятности предпочтения. Функция вознаграждения идентифицируется, таким образом, лишь с точностью до аддитивной константы.
\end{enumerate}
\end{proposition}

\begin{proof}
(1) следует из тождества $\sigma(z) + \sigma(-z) = 1$. (2) следует из $\sigma(0) = 1/2$. (3) следует из $\sigma'(z) = \sigma(z)(1-\sigma(z)) > 0$ для всех $z\in\R$. (4): подстановка $r^*(x,y) + c$ вместо $r^*(x,y)$ даёт $\sigma((r^*(x,y_1)+c) - (r^*(x,y_2)+c)) = \sigma(r^*(x,y_1)-r^*(x,y_2))$, что неизменно.
\end{proof}

\subsection{Связь с рейтингами Эло}

\index{Elo rating}
Модель Bradley--Terry тесно связана с \textbf{системой рейтингов Эло}, используемой в шахматах и других соревновательных играх. В Эло у каждого игрока $i$ есть скалярный рейтинг $\rho_i$, и вероятность того, что игрок~$i$ победит игрока~$j$, моделируется как $1/(1 + 10^{(\rho_j - \rho_i)/400})$. Это в точности модель Bradley--Terry с $r^* = \rho$ и логистической функцией по основанию~$10$ с масштабным коэффициентом~$400$. Ключевое различие: рейтинги Эло обновляются онлайн после каждого матча, используя простой шаг стохастического градиента по лог-потере, тогда как при моделировании вознаграждения мы выполняем пакетный градиентный спуск по всем парам предпочтений одновременно. Оба подхода оценивают одно и то же скрытое качество $r^*$ по попарным сравнениям.

% ===========================================================
\section{Модели вознаграждения по предпочтениям}
\label{sec:preference-rm}

\subsection{Оценка максимального правдоподобия}

Имея датасет предпочтений $\mathcal{D}$, мы хотим подобрать параметрическую модель $r_\psi$ для аппроксимации $r^*$. При предположении Bradley--Terry вероятность наблюдения полного датасета равна
\begin{equation}
\label{eq:bt-likelihood}
\mathcal{L}_{\text{lik}}(\psi)
= \prod_{i=1}^{N} \sigma\!\bigl(r_\psi(x^{(i)}, y_w^{(i)}) - r_\psi(x^{(i)}, y_l^{(i)})\bigr).
\end{equation}
Максимизация правдоподобия эквивалентна минимизации отрицательного логарифма правдоподобия.

\begin{definition}[Функция потерь модели вознаграждения по предпочтениям]
\label{def:rm-loss}
\index{reward model loss}
\textbf{Функция потерь модели вознаграждения по предпочтениям} равна
\begin{equation}
\label{eq:rm-loss}
\boxed{
\mathcal{L}_{\mathrm{RM}}(\psi)
= -\frac{1}{N}\sum_{i=1}^{N}
\log \sigma\!\bigl(r_\psi(x^{(i)}, y_w^{(i)}) - r_\psi(x^{(i)}, y_l^{(i)})\bigr).
}
\end{equation}
\end{definition}

Это в точности бинарная кросс-энтропийная потеря с меткой~$1$ («победитель предпочтителен») и предсказанной вероятностью $\hat{p}_i = \sigma(r_\psi(x^{(i)}, y_w^{(i)}) - r_\psi(x^{(i)}, y_l^{(i)}))$. Обучение модели вознаграждения --- это, следовательно, стандартная задача бинарной классификации: модель должна правильно ранжировать предпочтительный ответ выше непредпочтительного.

\begin{proposition}[Структура градиента $\mathcal{L}_{\mathrm{RM}}$]
\label{prop:rm-gradient}
Для одной пары $(x, y_w, y_l)$ пусть $\Delta = r_\psi(x,y_w) - r_\psi(x,y_l)$. Тогда
\[
\frac{\partial \mathcal{L}_{\mathrm{RM}}}{\partial r_\psi(x,y_w)} = -\sigma(-\Delta) < 0, \qquad
\frac{\partial \mathcal{L}_{\mathrm{RM}}}{\partial r_\psi(x,y_l)} = +\sigma(-\Delta) > 0.
\]
Градиентный спуск, таким образом, \emph{увеличивает} $r_\psi(x,y_w)$ и \emph{уменьшает} $r_\psi(x,y_l)$, раздвигая оценки. Размер шага $\sigma(-\Delta) = 1 - \sigma(\Delta)$ велик, когда текущий разрыв мал (пара сложная), и мал, когда разрыв уже велик (пара лёгкая). Это стандартное поведение логистической регрессии.
\end{proposition}

\begin{proof}
$\mathcal{L} = -\log\sigma(\Delta)$. Дифференцируя по $\Delta$: $\partial\mathcal{L}/\partial\Delta = -\sigma(-\Delta)$. Поскольку $\partial\Delta/\partial r_\psi(x,y_w) = +1$ и $\partial\Delta/\partial r_\psi(x,y_l) = -1$, результат следует по правилу цепочки.
\end{proof}

\subsection{Выбор архитектуры}

Модель вознаграждения по предпочтениям $r_\psi$ обычно строится поверх предобученной языковой модели. Каноническая архитектура инициализирует $r_\psi$ из контрольной точки SFT\,---\,той же модели с контролируемым дообучением, используемой в качестве отправной точки для RL. Линейный проекционный слой $W \in \R^{d \times 1}$ добавляется поверх трансформера, принимая на вход скрытое состояние на позиции последнего токена (обычно токена \texttt{[EOS]}) и выдавая скалярную оценку.

\begin{itemize}[leftmargin=*]
  \item \textbf{Почему инициализировать из SFT?} Модель вознаграждения должна оценивать ответы в том же формате, в котором модель SFT их генерирует. Инициализация из базовой LM, не видевшей данных для следования инструкциям, ведёт к модели вознаграждения, плохо понимающей содержание, которое ей следует оценивать.
  \item \textbf{Почему использовать токен EOS?} К тому времени, когда модель обрабатывает токен \texttt{[EOS]}, механизм само-внимания позволил каждой позиции обращать внимание на каждую другую. Скрытое состояние EOS, таким образом, кодирует сжатое представление всего ответа, делая его естественной точкой извлечения скалярной оценки качества.
  \item \textbf{Инвариантность к масштабу}: функция вознаграждения идентифицируется лишь с точностью до аддитивной константы (Утверждение~\ref{prop:bt-properties}). На практике оценки нормализуются до нулевого среднего и единичной дисперсии по обучающему батчу, что стабилизирует обучение и обеспечивает работу штрафа KL (Раздел~\ref{sec:overoptimization}) при согласованном масштабе.
\end{itemize}

\subsection{Практическая процедура обучения}

\begin{algorithm}[ht]
\caption{Обучение модели вознаграждения по предпочтениям}
\label{alg:rm-training}
\KwIn{Контрольная точка SFT $\pi_{\mathrm{SFT}}$; датасет предпочтений $\mathcal{D}$; скорость обучения $\alpha$; размер батча $B$; число эпох $E$}
\KwOut{Обученная модель вознаграждения $r_\psi$}
Инициализировать $r_\psi \leftarrow \pi_{\mathrm{SFT}}$ + случайная линейная голова $W \in \R^{d\times 1}$\;
\For{$e = 1, \ldots, E$}{
  \For{каждый мини-батч $\{(x^{(i)}, y_w^{(i)}, y_l^{(i)})\}_{i=1}^{B} \subset \mathcal{D}$}{
    Вычислить $s_w^{(i)} \leftarrow r_\psi(x^{(i)}, y_w^{(i)})$ и $s_l^{(i)} \leftarrow r_\psi(x^{(i)}, y_l^{(i)})$ для всех $i$\;
    Вычислить потерю $\mathcal{L} \leftarrow -\frac{1}{B}\sum_{i=1}^B \log \sigma(s_w^{(i)} - s_l^{(i)})$\;
    $\psi \leftarrow \psi - \alpha\,\nabla_\psi\,\mathcal{L}$\;
    Мониторировать $\mathrm{Acc} = \frac{1}{B}\sum_i \ind[s_w^{(i)} > s_l^{(i)}]$\;
  }
}
\Return $r_\psi$
\end{algorithm}

На практике $y_w$ и $y_l$ конкатенируются по батч-измерению, так что оба обрабатываются за один прямой проход, вдвое сокращая количество прямых проходов трансформера по сравнению с наивной двухпроходной реализацией. Хорошо обученные модели вознаграждения по предпочтениям обычно достигают попарной точности на отложенном наборе 70--80\%. Идеальная точность невозможна, поскольку человеческие аннотации содержат подлинный шум: разные аннотаторы расходятся во мнениях, и даже один аннотатор может быть непоследовательным в разных сессиях.

% ===========================================================
\section{Модели вознаграждения по исходам (ORM)}
\label{sec:orm}

\subsection{Разреженные терминальные вознаграждения}

Модель вознаграждения по предпочтениям из Раздела~\ref{sec:preference-rm} присваивает единственный скаляр полному ответу $(x, y)$. Это частный случай более широкого класса \textbf{моделей вознаграждения по исходам}\index{outcome reward model (ORM)} (ORM): моделей, оценивающих только конечный исход траектории, с нулевым вознаграждением на всех промежуточных шагах.

\begin{definition}[Модель вознаграждения по исходам]
\label{def:orm}
\index{outcome reward model (ORM)}
\textbf{Модель вознаграждения по исходам} (ORM) --- это функция вознаграждения $r_{\mathrm{ORM}} : \mathcal{X} \times \mathcal{Y} \to \R$, присваивающая скаляр полному ответу $y$, с нулевым вознаграждением на каждом промежуточном шаге генерации:
\[
R_t = 0 \mbox{ при } t < T, \qquad R_T = r_{\mathrm{ORM}}(x, y).
\]
\end{definition}

ORM --- естественный выбор, когда качество ответа можно оценить только после того, как полный ответ доступен. Они возникают как минимум в трёх различных контекстах.

\textbf{1. RLHF на основе предпочтений.} Как выведено в Разделе~\ref{sec:preference-rm}, результатом модели вознаграждения по предпочтениям на основе Bradley--Terry является единственный скаляр на ответ. Это ORM по конструкции.

\textbf{2. Задачи с верифицируемо правильным ответом.} При решении математических задач, программировании и фактическом вопросоответе ответ либо верен, либо нет: сгенерированный ответ совпадает с эталоном или нет. Можно обучить бинарный классификатор $r_{\mathrm{ORM}}(x, y) \in \{0, 1\}$, проверяющий, верен ли окончательный ответ. Это даёт чистый объективный ORM, избегающий шума, присущего аннотациям людей.

\textbf{3. Оценка ценности методом Монте-Карло.} В фреймворке Монте-Карло (MC) из Главы~\ref{ch:mc} оценки ценности получаются путём развёртки полных траекторий из каждого состояния и усреднения доходностей. ORM в контексте языковых моделей соответствует в точности доходности Монте-Карло: полная последовательность генерируется, терминальное вознаграждение наблюдается, а ценность каждого промежуточного состояния оценивается усреднением по многим развёрткам. Эта связь точна: сигнал вознаграждения ORM --- это доходность MC с нулевым промежуточным вознаграждением и терминальным вознаграждением.

\subsection{Когда использовать ORM}

ORM предпочтительны, когда:
\begin{itemize}[leftmargin=*]
  \item Корректность промежуточных шагов сложно или дорого оценить.
  \item Доступен надёжный автоматизированный верификатор (например, символьная программа проверки математики или запускатель юнит-тестов).
  \item Горизонт задачи достаточно короток, чтобы разреженное терминальное вознаграждение не приводило к неразрешимому присвоению кредита.
\end{itemize}

Главное ограничение ORM --- именно проблема присвоения кредита, подчёркнутая в Главе~\ref{ch:td}: ошибочный токен в начале ответа вносит вклад в неудачу всего ответа, но получает только косвенный сигнал через терминальное вознаграждение после множества промежуточных токенов. Модели вознаграждения по процессу, вводимые далее, устраняют это ограничение, предоставляя плотные промежуточные сигналы.

% ===========================================================
\section{Модели вознаграждения по процессу (PRM)}
\label{sec:prm}

\subsection{Надзор на уровне шагов}

Вместо оценки только конечного ответа \textbf{модель вознаграждения по процессу}\index{process reward model (PRM)} (PRM) присваивает вознаграждение каждому промежуточному шагу в генерации. В контексте многошагового математического рассуждения шаг соответствует строке решения (предложению, логическому выводу или алгебраической манипуляции). В более общем случае в RL языковых моделей шаг может соответствовать предложению, абзацу или даже отдельному токену.

\begin{definition}[Модель вознаграждения по процессу]
\label{def:prm}
\index{process reward model (PRM)}
\textbf{Модель вознаграждения по процессу} (PRM) присваивает скалярную оценку каждому шагу $k = 1, \ldots, K$ многошагового решения $y = (y_1, \ldots, y_K)$:
\[
r_{\mathrm{PRM}}(x, y_{1:k}) \in \R \quad \mbox{для каждого } k = 1, \ldots, K.
\]
Вознаграждение на шаге $k$ отражает качество или корректность частичного решения $(y_1, \ldots, y_k)$ при промпте $x$.
\end{definition}

\begin{figure}[ht]
\centering
\begin{tikzpicture}[
  node distance=0.4cm and 1.6cm,
  every node/.style={font=\small},
  stepbox/.style={draw, rounded corners=2pt, minimum width=1.5cm, minimum height=0.7cm, fill=blue!6},
  reward/.style={draw, circle, minimum size=0.55cm, inner sep=1pt, font=\tiny},
  arrow/.style={->, >=Stealth}
]
  % Строка ORM
  \node[font=\footnotesize\bfseries] (orm-label) at (-1.0, 1.0) {ORM};
  \node[stepbox] (o1) at (1.0, 1.0)  {Шаг 1};
  \node[stepbox] (o2) at (2.8, 1.0)  {Шаг 2};
  \node[stepbox] (o3) at (4.6, 1.0)  {Шаг 3};
  \node[stepbox] (o4) at (6.4, 1.0)  {Шаг 4};
  \node[reward, fill=green!20] (or4) at (8.0, 1.0) {$r$};
  \draw[arrow] (o1) -- (o2); \draw[arrow] (o2) -- (o3); \draw[arrow] (o3) -- (o4);
  \draw[arrow] (o4) -- (or4);
  \draw[dashed, gray] (o1.south) -- ++(0,-0.15) node[below, gray, font=\tiny] {$0$};
  \draw[dashed, gray] (o2.south) -- ++(0,-0.15) node[below, gray, font=\tiny] {$0$};
  \draw[dashed, gray] (o3.south) -- ++(0,-0.15) node[below, gray, font=\tiny] {$0$};

  % Строка PRM
  \node[font=\footnotesize\bfseries] (prm-label) at (-1.0,-1.0) {PRM};
  \node[stepbox] (p1) at (1.0,-1.0)  {Шаг 1};
  \node[stepbox] (p2) at (2.8,-1.0)  {Шаг 2};
  \node[stepbox] (p3) at (4.6,-1.0)  {Шаг 3};
  \node[stepbox] (p4) at (6.4,-1.0)  {Шаг 4};
  \node[reward, fill=green!20]  (pr1) at (1.0,-2.1)  {$r_1$};
  \node[reward, fill=yellow!30] (pr2) at (2.8,-2.1)  {$r_2$};
  \node[reward, fill=red!20]    (pr3) at (4.6,-2.1)  {$r_3$};
  \node[reward, fill=green!20]  (pr4) at (6.4,-2.1)  {$r_4$};
  \draw[arrow] (p1) -- (p2); \draw[arrow] (p2) -- (p3); \draw[arrow] (p3) -- (p4);
  \draw[arrow] (p1) -- (pr1); \draw[arrow] (p2) -- (pr2);
  \draw[arrow] (p3) -- (pr3); \draw[arrow] (p4) -- (pr4);
\end{tikzpicture}
\caption{Структура сигнала ORM и PRM. ORM присваивает единственное терминальное вознаграждение (вверху); PRM присваивает вознаграждение каждому шагу (внизу). Цвета шагов указывают качество: зелёный = верно, жёлтый = неопределённо, красный = неверно.}
\label{fig:orm-prm}
\end{figure}

\subsection{Обучение моделей вознаграждения по процессу}

Обучение PRM требует аннотаций на \emph{уровне шагов}: для каждого шага в решении аннотатор (человек или автоматизированный) указывает, является ли шаг верным. Пусть $c_k \in \{0, 1\}$ --- метка корректности для шага $k$. PRM обучается предсказывать вероятность того, что шаг $k$ верен:
\begin{equation}
\label{eq:prm-loss}
\mathcal{L}_{\mathrm{PRM}}(\psi) = -\frac{1}{N}\sum_{i=1}^{N}\frac{1}{K^{(i)}}\sum_{k=1}^{K^{(i)}} \Bigl[c_k^{(i)}\log r_\psi(x^{(i)}, y_{1:k}^{(i)}) + (1-c_k^{(i)})\log(1 - r_\psi(x^{(i)}, y_{1:k}^{(i)}))\Bigr],
\end{equation}
где $N$ --- число решений, а $K^{(i)}$ --- количество шагов в решении $i$.

\begin{example}[Датасет PRM800K]
\label{ex:prm800k}
\index{PRM800K}
Датасет PRM800K (Lightman et al., 2023) --- наиболее широко используемый бенчмарк для обучения и оценки PRM в математическом рассуждении. Он содержит около $800~000$ пошаговых человеческих аннотаций на решения задач из датасета MATH, коллекции олимпиадных математических задач семи уровней сложности. Каждый шаг каждого решения помечается как \emph{положительный} (шаг верен и продвигает решение), \emph{отрицательный} (шаг содержит ошибку) или \emph{нейтральный} (шаг нерелевантен или дублирован). PRM, обученная на PRM800K, учится присваивать высокую вероятность корректности логически верным шагам вывода и низкую вероятность ошибочным шагам, независимо от того, оказался ли конечный ответ верным.

Что важно, PRM800K демонстрирует, что надзор на уровне шагов даёт сигнал, недоступный ORM: решение может прийти к правильному окончательному ответу через неверный промежуточный шаг (случайные ошибки), или содержать все верные шаги, но прийти к неверному ответу из-за арифметической ошибки в последней строке. ORM присвоила бы высокое вознаграждение первому решению и низкое второму, несмотря на то что второе демонстрирует лучшее математическое рассуждение.
\end{example}

\subsection{Использование PRM для поиска во время вывода}

Помимо роли сигнала вознаграждения во время RL-обучения, PRM особенно мощны во время вывода. По заданному промпту $x$ можно сгенерировать несколько кандидатных решений с помощью лучевого поиска или best-of-$n$ сэмплирования, затем ранжировать их с помощью PRM. Несколько стратегий оказались эффективными:

\begin{description}[leftmargin=!,labelwidth=3.5cm]
  \item[Переранжирование best-of-$n$] Сгенерировать $n$ полных решений, оценить каждое с помощью ORM или агрегированных оценок PRM и вернуть решение с наивысшей оценкой. С хорошей PRM это существенно эффективнее, чем использование ORM для переранжирования, поскольку PRM штрафует решения, приходящие к верному ответу через ошибочное рассуждение.
  \item[Пошаговый лучевой поиск] На каждом шаге генерировать $K$ кандидатных продолжений, оценивать каждое с помощью PRM и расширять только $B$ лучших лучей. Это вычислительно дорого, но исследует пространство решений значительно эффективнее, чем плоское сэмплирование.
  \item[Древовидный поиск под руководством PRM] Использовать PRM как функцию ценности в древовидном поиске (например, Monte Carlo Tree Search) по шагам рассуждения. Это позволяет модели возвращаться назад и исследовать альтернативные пути рассуждения, когда ранний шаг получает низкую оценку PRM.
\end{description}

\subsection{ORM против PRM: сравнение}

\begin{table}[ht]
\centering
\caption{Сравнение моделей вознаграждения по исходам (ORM) и по процессу (PRM).}
\label{tab:orm-prm}
\small
\begin{tabular}{@{}lp{4.5cm}p{4.5cm}@{}}
\toprule
\textbf{Свойство} & \textbf{ORM} & \textbf{PRM} \\
\midrule
Гранулярность & Только исход & На уровне шагов \\
Стоимость аннотирования & Низкая & Высокая \\
Присвоение кредита & Разреженное & Плотное \\
Случайные ошибки & Вознаграждает любой верный ответ & Штрафует неверные шаги \\
Устойчивость к шагам & Низкая & Высокая \\
Качество поиска & Умеренное & Сильное \\
Необходимые данные & Результаты (верно/неверно) & Метки корректности шагов \\
Применения & RLHF по предпочтениям, верификация & Мат. рассуждение, код \\
\bottomrule
\end{tabular}
\end{table}

% ===========================================================
\section{Цели обучения и практические соображения}
\label{sec:rm-objectives}

\subsection{Варианты функции потерь с запасом}

Стандартная функция потерь модели вознаграждения~\eqref{eq:rm-loss} толкает $r_\psi(x, y_w)$ выше $r_\psi(x, y_l)$ без указания, насколько именно. Естественное расширение --- требовать минимальный запас $m > 0$:

\begin{equation}
\label{eq:rm-margin-loss}
\mathcal{L}_{\mathrm{RM}}^{(m)}(\psi)
= -\frac{1}{N}\sum_{i=1}^{N}
\log \sigma\!\bigl(r_\psi(x^{(i)}, y_w^{(i)}) - r_\psi(x^{(i)}, y_l^{(i)}) - m\bigr).
\end{equation}

Запас гарантирует, что правильно ранжированные пары с малым разрывом всё же вносят нетривиальный вклад в потерю, поощряя модель производить уверенно разделённые оценки. Это аналогично запасу в методах опорных векторов.

Связанный вариант ассоциирует непрерывный показатель уверенности $\delta^{(i)} \in [0, 1]$ с каждой парой предпочтений, отражающий уверенность аннотатора. Парам, с которыми аннотаторы сильно согласны, присваиваются более высокие $\delta$:
\begin{equation}
\label{eq:rm-weighted-loss}
\mathcal{L}_{\mathrm{RM}}^{(\delta)}(\psi)
= -\frac{1}{\sum_i \delta^{(i)}}\sum_{i=1}^{N}
\delta^{(i)}\,\log \sigma\!\bigl(r_\psi(x^{(i)}, y_w^{(i)}) - r_\psi(x^{(i)}, y_l^{(i)})\bigr).
\end{equation}
Это снижает вес неоднозначных пар предпочтений, уменьшая влияние шума в процессе аннотирования.

\subsection{Сглаживание меток}

Аннотации предпочтений людей шумны: аннотаторы иногда ошибаются, и некоторые пары подлинно неоднозначны. Сглаживание меток регуляризует модель вознаграждения, заменяя жёсткие метки $\{0,1\}$ мягкими метками $\{\epsilon/2, 1 - \epsilon/2\}$ для малого параметра сглаживания $\epsilon \in (0, 0.2]$:
\begin{equation}
\label{eq:rm-label-smooth}
\mathcal{L}_{\mathrm{RM}}^{(\epsilon)}(\psi)
= -\frac{1}{N}\sum_{i=1}^{N}
\Bigl[(1-\tfrac{\epsilon}{2})\log\sigma(\Delta^{(i)}) + \tfrac{\epsilon}{2}\log\sigma(-\Delta^{(i)})\Bigr],
\end{equation}
где $\Delta^{(i)} = r_\psi(x^{(i)}, y_w^{(i)}) - r_\psi(x^{(i)}, y_l^{(i)})$. Сглаживание меток не даёт модели вознаграждения гнать $\Delta^{(i)} \to +\infty$ (что сделало бы $\log\sigma(\Delta^{(i)}) \to 0$ и потерю нулевой) и побуждает к откалиброванным оценкам вероятности.

\subsection{Качество данных и протоколы аннотирования}

Качество модели вознаграждения ограничено качеством данных о предпочтениях. Ряд практических соображений существенен.

\begin{itemize}[leftmargin=*]
  \item \textbf{Согласие аннотаторов}: пары предпочтений, с которыми аннотаторы сильно расходятся, следует фильтровать или понижать в весе. Межаннотаторское согласие (например, каппа Коэна $\kappa > 0.4$) --- минимальная планка для надёжного надзора.
  \item \textbf{Разнообразие промптов}: модель вознаграждения, обученная на узком распределении промптов, плохо обобщается. Данные о предпочтениях должны охватывать разнообразные темы, форматы и намерения пользователей.
  \item \textbf{Разнообразие ответов}: если оба ответа $y_w$ и $y_l$ сэмплированы из сильной стратегии, разрыв предпочтения может быть мал, а сигнал --- шумным. Включение сравнений со значительно более слабыми базовыми линиями (например, с базовой моделью до SFT) может дать более лёгкий обучающий сигнал.
  \item \textbf{Оценка на отложенном наборе}: точность модели вознаграждения всегда следует сообщать на отложенном наборе из того же процесса аннотирования, и в идеале также на независимо собранном оценочном наборе для защиты от переобучения к идиосинкразиям аннотаторов.
\end{itemize}

\subsection{Метрики оценки}

Основная метрика оценки для модели вознаграждения по предпочтениям --- \textbf{попарная точность}\index{pairwise accuracy}:
\begin{equation}
\label{eq:rm-accuracy}
\mathrm{Acc} = \frac{1}{N_{\mathrm{eval}}}\sum_{i=1}^{N_{\mathrm{eval}}} \ind\!\bigl[r_\psi(x^{(i)}, y_w^{(i)}) > r_\psi(x^{(i)}, y_l^{(i)})\bigr].
\end{equation}

Помимо попарной точности можно вычислять \textbf{корреляцию Пирсона} между оценками модели вознаграждения и независимыми человеческими рейтингами, и \textbf{метрики ранжирования}, такие как $\tau$ Кенделла или $\rho$ Спирмена, по спискам ответов, ранжированных по оценкам модели вознаграждения. Для ORM, используемых в верифицируемых задачах, точность относительно символьного верификатора служит прямой проверкой относительно истины.

% ===========================================================
\section{Переоптимизация вознаграждения}
\label{sec:overoptimization}

\subsection{Закон Гудхарта в обучении с подкреплением}

\index{Goodhart's Law}
\index{reward overoptimisation}
Обученная модель вознаграждения $r_\psi$ является прокси для истинного человеческого предпочтения $r^*$. Когда стратегия оптимизируется против $r_\psi$ с помощью RL, она неизбежно эксплуатирует любое расхождение между $r_\psi$ и $r^*$. Это частный случай \textbf{закона Гудхарта}: \emph{когда мера становится целью, она перестаёт быть хорошей мерой.}

\begin{definition}[Переоптимизация вознаграждения]
\label{def:overoptimization}
\index{reward overoptimisation}
\textbf{Переоптимизация вознаграждения} (или \emph{взлом вознаграждения} в смысле прокси) происходит, когда стратегия $\pi_\theta$ достигает высокого прокси-вознаграждения $\E[r_\psi(x,y)]$, но низкого истинного вознаграждения $\E[r^*(x,y)]$. Она характеризуется существованием KL-бюджета $d^*$ такого, что:
\begin{align}
\frac{d}{d\,d}\E[r^*(x,y)] &> 0 \quad \mbox{при } d < d^*, \\
\frac{d}{d\,d}\E[r^*(x,y)] &< 0 \quad \mbox{при } d > d^*,
\end{align}
где $d = \KL{\pi_\theta}{\pi_{\mathrm{ref}}}$ измеряет дивергенцию от опорной стратегии.
\end{definition}

Иными словами: по мере того, как стратегия всё дальше оптимизируется против прокси-вознаграждения (с возрастающей дивергенцией KL от опорной), истинное качество сначала улучшается, достигает пика при $d^*$, а затем деградирует. Прокси-вознаграждение, тем временем, продолжает монотонно улучшаться. Это создаёт опасный разрыв: стандартные RL-метрики обучения выглядят хорошо, даже когда модель ухудшается по той мере, что действительно важна.

\begin{example}[Масштабирование переоптимизации]
\label{ex:overopt-scaling}
Gao et al.\ (2023) охарактеризовали это явление эмпирически и обнаружили, что истинное вознаграждение $r^*$ как функция дивергенции KL $d$ следует приблизительно вогнутой кривой вида
\[
r^*(d) \approx \alpha\,\sqrt{d} - \beta\,d,
\]
для положительных констант $\alpha$ и $\beta$, зависящих от размера и качества модели вознаграждения. Оптимальный KL-бюджет равен $d^* = (\alpha/(2\beta))^2$. Что важно, $\beta$ растёт по мере того, как модель вознаграждения становится \emph{меньше} относительно оптимизируемой стратегии: слабая модель вознаграждения ведёт к более быстрой переоптимизации. Это говорит о том, что ёмкость модели вознаграждения должна масштабироваться с ёмкостью стратегии.
\end{example}

\subsection{Штраф KL как регуляризация}

Стандартное средство от переоптимизации вознаграждения --- добавить штраф дивергенции KL к целевой функции, ограничивая стратегию оставаться близко к опорной (SFT) стратегии:

\begin{definition}[KL-регуляризованная целевая функция RLHF]
\label{def:kl-rlhf}
\index{KL regularisation}
\textbf{KL-регуляризованная целевая функция RLHF} равна
\begin{equation}
\label{eq:kl-rlhf}
J_{\mathrm{RLHF}}(\theta)
= \E_{x \sim \mathcal{D},\; y \sim \pi_\theta(\cdot|x)}\!\bigl[r_\psi(x, y)\bigr]
- \beta\,\KL{\pi_\theta(\cdot|x)}{\pi_{\mathrm{ref}}(\cdot|x)},
\end{equation}
где $\beta > 0$ --- коэффициент KL, а $\pi_{\mathrm{ref}}$ --- стратегия SFT.
\end{definition}

В реализациях PPO (Глава~\ref{ch:practical-rlhf}) штраф KL обычно включается в сигнал вознаграждения на уровне токенов:
\begin{equation}
\label{eq:kl-token-reward}
r(x, y) = r_\psi(x, y) - \beta\,\log\frac{\pi_\theta(y \mid x)}{\pi_{\mathrm{ref}}(y \mid x)},
\end{equation}
что эквивалентно добавлению $-\beta\log(\pi_\theta/\pi_{\mathrm{ref}})$ как бонуса формирования на уровне токенов. Это в точности функция формирования на основе потенциала с $\Phi(s) = -\beta\log\pi_{\mathrm{ref}}(a|s)$\,---\,связывая штраф KL с потенциальным фреймворком Раздела~\ref{sec:reward-shaping}.

\begin{theorem}[Оптимальная стратегия при KL-регуляризации]
\label{thm:kl-optimal-policy}
\index{KL regularisation!optimal policy}
Стратегия, максимизирующая~\eqref{eq:kl-rlhf} поточечно (для каждого $x$), равна
\begin{equation}
\label{eq:kl-optimal-policy}
\pi^*(y \mid x) = \frac{\pi_{\mathrm{ref}}(y \mid x)\,\exp\!\bigl(r_\psi(x, y)/\beta\bigr)}{Z(x)},
\end{equation}
где $Z(x) = \sum_y \pi_{\mathrm{ref}}(y \mid x)\,\exp(r_\psi(x, y)/\beta)$ --- статистическая сумма.
\end{theorem}

\begin{proof}
Зафиксируем $x$ и запишем целевую функцию как функционал от $\pi(\cdot|x)$:
\[
\sum_y \pi(y|x)\,r_\psi(x,y) - \beta\sum_y \pi(y|x)\log\frac{\pi(y|x)}{\pi_{\mathrm{ref}}(y|x)}.
\]
Объединяя слагаемые:
\[
= -\beta\sum_y \pi(y|x)\left[\log\frac{\pi(y|x)}{\pi_{\mathrm{ref}}(y|x)} - \frac{r_\psi(x,y)}{\beta}\right]
= -\beta\,\KL{\pi(\cdot|x)}{\tilde\pi(\cdot|x)} + \text{const},
\]
где $\tilde\pi(y|x) \propto \pi_{\mathrm{ref}}(y|x)\exp(r_\psi(x,y)/\beta)$. KL минимизируется (и целевая функция максимизируется) при $\pi = \tilde\pi$, давая~\eqref{eq:kl-optimal-policy}.
\end{proof}

Формула~\eqref{eq:kl-optimal-policy} раскрывает роль $\beta$ как параметра \emph{температуры}:
\begin{itemize}[leftmargin=*]
  \item При $\beta \to 0$: $\pi^*$ концентрирует всю массу на ответе, максимизирующем $r_\psi$. Стратегия полностью управляется прокси-вознаграждением, и взлом вознаграждения максимален.
  \item При $\beta \to \infty$: $\pi^* \to \pi_{\mathrm{ref}}$. RL-дообучение не имеет эффекта; стратегия остаётся в контрольной точке SFT.
  \item При умеренных $\beta \in [0.01, 0.5]$: стратегия интерполирует между априором SFT и решением, максимизирующим вознаграждение. Эмпирически этот диапазон захватывает пик кривой истинного вознаграждения (Определение~\ref{def:overoptimization}).
\end{itemize}

\begin{remark}[Связь с DPO]
\label{rem:dpo-connection}
Формулу оптимальной стратегии~\eqref{eq:kl-optimal-policy} можно преобразовать, чтобы выразить $r_\psi$ через $\pi^*$, $\pi_{\mathrm{ref}}$ и $Z(x)$:
\[
r_\psi(x, y) = \beta\log\frac{\pi^*(y|x)}{\pi_{\mathrm{ref}}(y|x)} + \beta\log Z(x).
\]
Это тождество лежит в основе алгоритма Direct Preference Optimisation (DPO): подстановка этого выражения для $r_\psi$ в модель предпочтений Bradley--Terry~\eqref{eq:bt-model} и использование того факта, что $\log Z(x)$ сокращается в разности, позволяет DPO обойти модель вознаграждения полностью и оптимизировать стратегию непосредственно по данным предпочтений. DPO выводится подробно в Главе~\ref{ch:practical-rlhf}.
\end{remark}

% ===========================================================
\section{Резюме и связи}
\label{sec:ch10-summary}

В этой главе были разработаны два взаимодополняющих фреймворка для проблемы проектирования вознаграждения.

\textbf{Формирование вознаграждения на основе потенциала} (Раздел~\ref{sec:reward-shaping}) предоставляет теоретически обоснованный метод дополнения данной функции вознаграждения для более плотного сигнала обучения без изменения множества оптимальных стратегий. Ключевое наблюдение --- телескопическое тождество~\eqref{eq:shaping-return}: бонус на основе потенциала $F = \gamma\Phi(s') - \Phi(s)$ вносит ровно $-\Phi(s_0)$ в кумулятивную доходность, зависящую от состояния константу, которая не влияет на функции преимущества. Этот принцип лежит в основе формулировки преимущества в методах актор-критик и штрафа KL на уровне токенов в RLHF.

\textbf{Обученные модели вознаграждения} (Разделы~\ref{sec:reward-hacking}--\ref{sec:overoptimization}) рассматривают более сложную проблему вывода вознаграждения из человеческой обратной связи. Модель Bradley--Terry предоставляет принципиальный вероятностный фреймворк для данных предпочтений, приводящий к кросс-энтропийной функции потерь модели вознаграждения~\eqref{eq:rm-loss}. Модели вознаграждения по исходам (ORM) предоставляют единственный терминальный сигнал, аналогичный доходностям Монте-Карло; модели вознаграждения по процессу (PRM) предоставляют плотный надзор на уровне шагов, разрешающий неоднозначности, недоступные ORM. Обе подчиняются закону Гудхарта: продолжительная оптимизация против любого прокси-вознаграждения в конечном счёте деградирует истинное качество. KL-регуляризованная целевая функция~\eqref{eq:kl-rlhf} и её оптимальное решение~\eqref{eq:kl-optimal-policy} дают принципиальное средство.

Центральные связи с последующими главами следующие.

\begin{description}[leftmargin=!,labelwidth=3.5cm]
  \item[Глава~\ref{ch:practical-rlhf} (PPO)] Обученная здесь модель вознаграждения $r_\psi$ служит непосредственно сигналом вознаграждения при дообучении PPO. Вознаграждение на уровне токенов~\eqref{eq:kl-token-reward} --- это в точности вознаграждение, вычисляемое на каждом шаге в оценщике преимущества PPO.
  \item[Глава~\ref{ch:practical-rlhf} (DPO)] Параметризация оптимальной стратегии~\eqref{eq:kl-optimal-policy} непосредственно приводит к целевой функции DPO, устраняющей необходимость обучать отдельную модель вознаграждения.
  \item[Глава~\ref{ch:practical-rlhf} (GRPO)] Групповая относительная оптимизация стратегии использует ORM на основе верификатора в качестве сигнала вознаграждения и нормализует вознаграждения внутри группы выбранных ответов, используя групповую статистику для оценки преимущества без обученного критика.
\end{description}

\begin{remark}[Историческая заметка]
\label{rem:history}
Формирование вознаграждения на основе потенциала было введено Нгом, Харадой и Расселлом на ICML 1999 и быстро стало стандартным инструментом в арсенале практиков RL. Модель Bradley--Terry восходит к 1952 году и была введена в статистике задолго до её принятия в машинном обучении. Её применение к RLHF было популяризировано Christiano et al.\ (2017), показавшими, что модель вознаграждения по предпочтениям, обученная на небольшом числе человеческих сравнений, может направлять RL в задачах Atari и робототехники. Различие между ORM и PRM было кристаллизовано Lightman et al.\ (2023) в контексте математического рассуждения, причём датасет PRM800K предоставил первый крупномасштабный ресурс пошаговых аннотаций.
\end{remark}

% ===========================================================
\section*{Упражнения}
\addcontentsline{toc}{section}{Упражнения}

\begin{exercise}
\label{ex:shaping-1}
Пусть $\cS = \{s_1, s_2, s_3\}$ --- пространство состояний детерминированного MDP с $\gamma = 0.9$ и истинным вознаграждением $R(s_1 \to s_2) = 0$, $R(s_2 \to s_3) = 0$, $R(s_3) = +1$ (терминал). Определим потенциальную функцию $\Phi(s_1) = 0$, $\Phi(s_2) = 0.5$, $\Phi(s_3) = 1.0$.
\begin{enumerate}
  \item Вычислите бонус формирования $F(s_1, s_2)$ и $F(s_2, s_3)$.
  \item Вычислите сформированную доходность $\tilde{G}_0$, начиная из $s_1$, и проверьте, что $\tilde{G}_0 = G_0 - \Phi(s_1)$.
  \item Имеет ли сформированный MDP ту же оптимальную стратегию, что и исходный? Объясните, используя Теорему~\ref{thm:shaping}.
\end{enumerate}
\end{exercise}

\begin{exercise}
\label{ex:shaping-non-potential}
Постройте явный MDP и не-потенциальную функцию формирования $F(s,a,s')$, меняющую оптимальную стратегию. Ваш MDP должен иметь как минимум два действия хотя бы в одном состоянии, и вы должны идентифицировать, какое действие было оптимальным до формирования и какое становится оптимальным после. Проверьте, что $F$ не может быть записана в виде $\gamma\Phi(s') - \Phi(s)$ для какой-либо $\Phi$.
\end{exercise}

\begin{exercise}
\label{ex:bt-properties}
Докажите все четыре свойства модели Bradley--Terry, перечисленные в Утверждении~\ref{prop:bt-properties}. Для свойства~4 (инвариантность к сдвигу) объясните, почему это означает, что функция вознаграждения не является единственно идентифицируемой по данным предпочтений, и опишите дополнительное ограничение, которое можно наложить для её идентификации.
\end{exercise}

\begin{exercise}
\label{ex:rm-loss-worked}
Предположим, что модель вознаграждения в данный момент присваивает оценки $r_\psi(x, y_w) = 1.2$ и $r_\psi(x, y_l) = 1.8$ для одной пары предпочтений.
\begin{enumerate}
  \item Вычислите потерю $\mathcal{L}_{\mathrm{RM}}$ для этой пары.
  \item Вычислите градиенты $\partial\mathcal{L}/\partial r_\psi(x,y_w)$ и $\partial\mathcal{L}/\partial r_\psi(x,y_l)$ и опишите направление обновления.
  \item Как изменяется величина градиента при росте разрыва $r_\psi(x,y_w) - r_\psi(x,y_l)$ от $-0.6$ до $+5$? Набросайте градиент как функцию разрыва.
  \item Теперь предположим, что мы добавляем запас $m = 1.0$. Пересчитайте потерю и градиент. Увеличивается или уменьшается градиент? Объясните интуитивно.
\end{enumerate}
\end{exercise}

\begin{exercise}
\label{ex:orm-prm}
Языковую модель просят решить следующую арифметическую задачу: «Если поезд едет со скоростью 60 миль/ч в течение 2.5 часа, а затем со скоростью 80 миль/ч в течение 1.5 часа, каково общее расстояние?» Модель генерирует следующее трёхшаговое решение:
\begin{enumerate}
  \item[Шаг 1:] Расстояние на первом отрезке: $60 \times 2.5 = 150$ миль.
  \item[Шаг 2:] Расстояние на втором отрезке: $80 \times 2 = 160$ миль. [\emph{Ошибка: должно быть $1.5$ часа}]
  \item[Шаг 3:] Итого: $150 + 160 = 310$ миль. [\emph{Ошибка перенесена из Шага 2; верный ответ 270}]
\end{enumerate}
\begin{enumerate}
  \item Какое вознаграждение присвоит ORM, если верификатор проверяет только конечный ответ? Что присвоит PRM?
  \item Чем ORM и PRM отличаются в обращении с этим решением по сравнению с решением, приходящим к 270 по другому (но верному) трёхшаговому пути?
  \item Опишите сценарий обучения, при котором ORM систематически переоценивала бы качество модели, делающей систематические ошибки на Шаге 2.
\end{enumerate}
\end{exercise}

\begin{exercise}
\label{ex:kl-optimal-policy}
Используйте Теорему~\ref{thm:kl-optimal-policy} для ответа на следующие вопросы об оптимальной стратегии $\pi^*$ при KL-регуляризации.
\begin{enumerate}
  \item Покажите, что $\log\pi^*(y|x) - \log\pi_{\mathrm{ref}}(y|x) = r_\psi(x,y)/\beta - \log Z(x)$.
  \item Используя результат (а), выразите $r_\psi(x,y)$ полностью через $\pi^*$, $\pi_{\mathrm{ref}}$ и $Z(x)$. Это тождество, используемое в выводе DPO.
  \item Проверьте, что подстановка $r_\psi(x,y) - r_\psi(x,y')$ в формулу Bradley--Terry~\eqref{eq:bt-model} приводит к сокращению $\log Z(x)$, оставляя вероятность предпочтения в терминах только $\pi^*$ и $\pi_{\mathrm{ref}}$.
\end{enumerate}
\end{exercise}

\begin{exercise}
\label{ex:goodhart}
Сформулируйте явление переоптимизации вознаграждения как задачу смещения-дисперсии в следующем смысле: прокси-модель вознаграждения $r_\psi$ является оценкой истинного вознаграждения $r^*$ с некоторой ошибкой обобщения $\varepsilon = \E[(r_\psi(x,y) - r^*(x,y))^2]$.
\begin{enumerate}
  \item Аргументируйте, что разрыв $\E_{\pi_\theta}[r_\psi] - \E_{\pi_\theta}[r^*]$ (разрыв между прокси- и истинным вознаграждением при стратегии $\pi_\theta$) растёт с дивергенцией KL $\KL{\pi_\theta}{\pi_{\mathrm{ref}}}$.
  \item Предположим, что ошибка обобщения модели вознаграждения $\varepsilon$ уменьшается по мере сбора большего объёма данных предпочтений. Как это влияет на оптимальный KL-бюджет $d^*$?
  \item Опишите практический критерий остановки дообучения PPO, использующий как прокси-вознаграждение, так и отложенный набор человеческих оценок для обнаружения момента, когда разрыв между прокси- и истинным вознаграждением становится недопустимо большим.
\end{enumerate}
\end{exercise}

\begin{exercise}
\label{ex:label-smoothing}
Начиная со сглаженной по меткам функции потерь~\eqref{eq:rm-label-smooth}, покажите, что она эквивалентна стандартной функции потерь~\eqref{eq:rm-loss}, возмущённой слагаемым $\frac{\epsilon}{2}\log\sigma(-\Delta)/(1-\epsilon/2)$, и интерпретируйте это возмущение как регуляризатор, предотвращающий произвольный рост оценок вознаграждения.
\end{exercise}

\begin{exercise}
\label{ex:potential-kl}
Покажите, что штраф KL на уровне токенов в~\eqref{eq:kl-token-reward} можно интерпретировать как бонус формирования на основе потенциала в MDP уровня токенов, обсуждаемом в разделе~\ref{sec:reward-hacking}. В частности:
\begin{enumerate}
  \item Определите потенциал $\Phi(s_t) = \beta\log\pi_{\mathrm{ref}}(a_t | s_t)$, вычисленный в каждой паре состояние-действие MDP.
  \item Вычислите бонус формирования $F(s_t, a_t, s_{t+1}) = \gamma\Phi(s_{t+1}) - \Phi(s_t)$ и сравните его со слагаемым KL на уровне токенов.
  \item Заключите, что KL-регуляризованная целевая функция RLHF эквивалентна нерегуляризованной RL-задаче на сформированном MDP, и укажите, какое следствие Теорема~\ref{thm:shaping} имеет для этой эквивалентности.
\end{enumerate}
\end{exercise}
